\documentclass[10pt,a4paper]{article}
\usepackage[landscape,margin=12mm,top=19mm,bottom=15mm,headheight=25pt,footskip=8mm]{geometry}
\usepackage[utf8]{inputenc}
\usepackage[T1]{fontenc}
\usepackage[brazil]{babel}
\usepackage{helvet,amsmath,amssymb,booktabs,array,tabularx,enumitem,graphicx,xcolor,tcolorbox,fancyhdr,hyperref,tikz}
\usetikzlibrary{arrows.meta,positioning}
\renewcommand{\familydefault}{\sfdefault}
\definecolor{azul}{RGB}{0,83,134}
\definecolor{ouro}{RGB}{206,177,68}
\hypersetup{pdflang=pt-BR,colorlinks=true,urlcolor=azul,linkcolor=azul,pdfauthor={Luan F. Nogueira e Luiz F. de Melo Soares}}
\setlength{\parindent}{0pt}
\setlength{\parskip}{5pt}
\setlist{nosep,leftmargin=*}
\pagestyle{fancy}\fancyhf{}
\fancyhead[L]{\small\color{azul}DO TRANSISTOR AO CHIP}
\fancyhead[R]{\includegraphics[height=7mm]{../slides/libs/emblemufc.pdf}}
\fancyfoot[L]{\scriptsize Luan F. Nogueira • Luiz F. de Melo Soares}
\fancyfoot[R]{\scriptsize\thepage}
\renewcommand{\headrulewidth}{.4pt}
\setlength{\headheight}{25pt}
\newcommand{\titulo}[2]{{\LARGE\color{azul}\bfseries #1}\par{\large #2}\par\vspace{2mm}}
\newcommand{\caixa}[2]{\begin{tcolorbox}[colback=azul!4,colframe=azul,title={#1},fonttitle=\bfseries,boxrule=.5pt,arc=1mm]#2\end{tcolorbox}}
\newcommand{\VDD}{V_{\mathrm{DD}}}

\hypersetup{pdftitle={Roadmap de estudos em microeletrônica}}
\begin{document}
\titulo{Para onde seguir na microeletrônica?}{Mapa de estudos · encerramento da aula 3}
\caixa{Uma base compartilhada}{Circuitos elétricos $\rightarrow$ semicondutores e MOSFET $\rightarrow$ CMOS e portas lógicas $\rightarrow$ simulação e análise de resultados. Acrescente matemática, sinais e sistemas, programação e documentação ao longo do percurso.}
\begingroup
\tcbset{height=8cm,before upper=\raggedright}
\begin{minipage}[t]{.19\linewidth}\vspace{0pt}
\caixa{Digital}{\textbf{Objetivo}\\Implementar funções em hardware.\par
\textbf{Caminho}\\Lógica combinacional\\$\downarrow$\\Registradores e estados\\$\downarrow$\\HDL/RTL e verificação\\$\downarrow$\\Síntese e análise temporal\par
\textbf{Primeiro projeto}\\ULA de 4 bits com testes automáticos.}
\end{minipage}\hfill
\begin{minipage}[t]{.19\linewidth}\vspace{0pt}
\caixa{Analógica}{\textbf{Objetivo}\\Processar sinais contínuos.\par
\textbf{Caminho}\\Polarização do MOSFET\\$\downarrow$\\Pequenos sinais\\$\downarrow$\\Realimentação e ruído\\$\downarrow$\\Variações e layout\par
\textbf{Primeiro projeto}\\Common-source: ganho e excursão de saída.}
\end{minipage}\hfill
\begin{minipage}[t]{.19\linewidth}\vspace{0pt}
\caixa{Sinais mistos}{\textbf{Objetivo}\\Conectar analógico e digital.\par
\textbf{Caminho}\\Base analógica e digital\\$\downarrow$\\Amostragem e quantização\\$\downarrow$\\ADCs e DACs\\$\downarrow$\\Integração e verificação\par
\textbf{Primeiro projeto}\\Modelo de ADC de 3 bits e erro de quantização.}
\end{minipage}\hfill
\begin{minipage}[t]{.19\linewidth}\vspace{0pt}
\caixa{RF integrada}{\textbf{Objetivo}\\Transmitir e receber sinais.\par
\textbf{Caminho}\\Analógica e sinais\\$\downarrow$\\Eletromagnetismo\\$\downarrow$\\Impedâncias e parâmetros S\\$\downarrow$\\Ruído e circuitos de RF\par
\textbf{Primeiro projeto}\\Rede LC de casamento: frequência e perdas.}
\end{minipage}\hfill
\begin{minipage}[t]{.19\linewidth}\vspace{0pt}
\caixa{Arquitetura}{\textbf{Objetivo}\\Organizar a computação.\par
\textbf{Caminho}\\ULA e registradores\\$\downarrow$\\Instruções e processador\\$\downarrow$\\Pipeline e memória\\$\downarrow$\\Paralelismo e aceleradores\par
\textbf{Primeiro projeto}\\Caminho de dados simples: medir ciclos de execução.}
\end{minipage}\par
\endgroup
\vspace{1mm}
\textbf{Outras conexões:} verificação e validação • projeto físico e layout • dispositivos e processos • teste e confiabilidade • encapsulamento e integração • ferramentas de projeto (EDA).

\colorbox{ouro!20}{\parbox{\dimexpr\linewidth-2\fboxsep}{\textbf{Próximo passo:} escolha uma pergunta $\rightarrow$ revise os pré-requisitos $\rightarrow$ faça um projeto pequeno $\rightarrow$ documente testes, resultados e limites. As trilhas se cruzam; o percurso não precisa ser linear.}}

{\footnotesize\textbf{Siglas:} ULA: unidade lógica e aritmética; HDL: linguagem de descrição de hardware; RTL: nível de transferências entre registradores; ADC/DAC: conversão analógico--digital/digital--analógica; RF: radiofrequência; LC: indutor e capacitor. Pipeline: execução organizada em etapas.\par
\textbf{Referências:} Salem e Dutertre (PDFs fornecidos); \href{https://ocw.mit.edu/courses/6-004-computation-structures-spring-2017/}{MIT 6.004}; visões de áreas por \href{https://www.cadence.com/en_US/home/explore/what-is-ic-design.html}{Cadence} e \href{https://www.synopsys.com/glossary/what-is-mixed-signal-integrated-circuit.html}{Synopsys}. Sequência e projetos: proposta didática dos facilitadores.}
\end{document}
